\section{Conclusion and future work}
We presented HCMO, the first ontology to model Home-Cage Monitoring as an integrated system spanning the animal subject, its housing, the environment, and the technical acquisition chain. HCMO captures HCM data together with the experimental context that makes them interpretable, keeps sensor, observation, and result distinct, and reuses selected terms from BFO, IAO, SOSA/SSN, Schema.org, OWL-Time, and QUDT rather than duplicating them. HCMO pins SOSA to the 2017 Recommendation and selectively reuses canonical SemTS 1.2.0 segment terms. A pinned instance-level ISA/Bioschemas, PROV-O, OBI, and STATO evidence slice demonstrates one acyclic recording-to-statistical-result workflow without claiming HCMO class mappings or formal ISA conformance. A generated 2 $\times$ 2 round-trip fixture further preserves animal identity, housing history, factors, groups, repeated observations, and separate semantic statistical results and their representing file fragments in HCMO RDF and an extended ISA RO-Crate. A narrower native ISA-JSON/ISA-Tab test preserves the real animal-Source-to-tissue-Sample collection without fabricating an animal Sample; all excluded HCMO/STATO structures are reported as controlled losses. HCMO is delivered as a FAIR, tool-consumable resource --- modular sources, reproducible distributions, SHACL shapes, competency-question queries, a JSON-LD context, persistent identifiers, an open licence, and documentation --- and is grounded in the COST TEATIME community.

\textbf{Limitations.} HCMO is still an early and evolving resource. Some conceptual boundaries are being stabilised, most notably between raw data, processed measurement, interpreted observation, and biological result, where the same information can be viewed at several levels. The diversity of HCM systems makes a single model that is both broad enough to be reusable and precise enough to describe individual experiments an ongoing balance. Alignment to external vocabularies is deliberately partial in this version. The factorial and mixed-model resources are interoperability fixtures, not a general HCMO statistics hierarchy. Complex pathologies and AI models on the data remain outside the stabilised core.

\textbf{Future work.} Planned directions include expert review of definitions and axioms identified by the class and property audits; systematic alignment decisions (reuse as-is vs specialise vs keep HCMO-specific) and bridge modules to adjacent ontologies (e.g. OBI, PROV-O, STATO, MEDO) and to taxa/anatomy/device vocabularies; extending the current selective QUDT pattern to additional quantity roles where justified; populating the ontology with real datasets via the TEATIME network and publishing the corresponding SPARQL queries; and supporting the downstream uses outlined in \S{}7. Through these steps, HCMO aims to become a community reference that improves the annotation, comparison, and reuse of HCM data across laboratories.
